\documentclass[a4paper,14pt]{extarticle}

\usepackage[left=30mm, right=10mm, top=20mm, bottom=20mm]{geometry}
\usepackage[T2A]{fontenc}
\usepackage[utf8]{inputenc}
\usepackage[russian]{babel}
\usepackage{tempora}
\usepackage{amsmath}
\usepackage{setspace}
\usepackage{enumitem}

\onehalfspacing

\begin{document}

\begin{center}
    \textbf{\Large Речь к защите курсовой работы (послайдовая)} \\
    \textit{Хронометраж: 3.5--4 минуты}
\end{center}

\vspace{1em}

\noindent \textbf{[Слайд 1: Титульный]} \\
Здравствуйте! Тема работы: \textbf{«Решение уравнений матфизики методом PINN и анализ устойчивости через NTK»}. Работа посвящена исследованию современных нейросетевых методов в задачах математической физики.

\vspace{0.8em}

\noindent \textbf{[Слайд 2: Постановка задачи]} \\
Объектом исследования выбрано уравнение Курамото-Сивашинского. Оно моделирует сложные волновые процессы и характеризуется балансом между неустойчивостью и стабилизацией. Наличие аналитического решения позволяет нам точно верифицировать результаты.

\vspace{0.8em}

\noindent \textbf{[Слайд 3: Метод PINN]} \\
Мы используем нейросети, информированные физикой. Главное отличие от классического обучения — включение невязки дифференциального уравнения непосредственно в функцию потерь, что гарантирует соблюдение физических законов.

\vspace{0.8em}

\noindent \textbf{[Слайд 4: Обучающая выборка]} \\
Для обучения формируется набор из начальных, граничных и внутренних коллокационных точек. Автоматическое дифференцирование позволяет вычислять производные произвольного порядка без использования сеточных методов.

\vspace{0.8em}

\noindent \textbf{[Слайд 5: Начало обучения]} \\
На первых этапах сеть восстанавливает лишь общую форму решения. Из-за наличия производной четвертого порядка модель проявляет склонность к низкочастотной аппроксимации на ранних итерациях.

\vspace{0.8em}

\noindent \textbf{[Слайд 6: Финальная аппроксимация]} \\
После завершения обучения базовая модель идеально восстанавливает профиль бегущей волны. На графиках видно полное совпадение нейросетевого решения с аналитическим эталоном.

\vspace{0.8em}

\noindent \textbf{[Слайд 7: Динамика ошибок]} \\
Суммарная ошибка достигла уровня $10^{-5}$. Компоненты функции потерь убывают стабильно, что подтверждает корректность настройки гиперпараметров базовой модели.

\vspace{0.8em}

\noindent \textbf{[Слайд 8: Гипотеза понижения порядка]} \\
Второй этап работы — проверка гипотезы: упростит ли обучение замена уравнения 4-го порядка на систему уравнений 2-го порядка через введение вспомогательной переменной $v = u_{xx}$.

\vspace{0.8em}

\noindent \textbf{[Слайд 9: Исследуемые архитектуры]} \\
Мы протестировали три подхода: единую сеть с двумя выходами, две независимые модели и гибридную архитектуру с общим стволом — \texttt{shared\_trunk}.

\vspace{0.8em}

\noindent \textbf{[Слайд 10: Сравнение точности]} \\
Гипотеза не подтвердилась: точность упала до $10^{-4}$. На графиках виден выход на «плато», связанный с трудностью одновременной минимизации невязок двух разных полей.

\vspace{0.8em}

\noindent \textbf{[Слайд 11: Анализ полей u и v]} \\
Поле $u$ восстанавливается корректно, однако вспомогательное поле $v$ демонстрирует сглаживание пиков. Это подтверждает, что прямое решение уравнения высокого порядка эффективнее декомпозиции.

\vspace{0.8em}

\noindent \textbf{[Слайд 12: Причины деградации]} \\
Основная причина — сложность синхронизации полей разной гладкости. Математическое упрощение лишь перенесло проблему в плоскость оптимизации и междмодельной связи.

\vspace{0.8em}

\noindent \textbf{[Слайд 13: NTK - Спектр собственных значений]} \\
Перейдем к диагностике через NTK. Рост ведущих собственных значений подтверждает \textbf{ускоренную сходимость} и устойчивость базовой модели к статистическим шумам.

\vspace{0.8em}

\noindent \textbf{[Слайд 14: NTK - Тепловая карта]} \\
Отсутствие разреженных зон на матрице взаимодействий гарантирует связность вычислений. Это обеспечивает глобальную гладкость и корректную передачу условий от границ вглубь области.

\vspace{0.8em}

\noindent \textbf{[Слайд 15: Анализ градиентов]} \\
Градиентный поток распределен равномерно. Отсутствие затухания подтверждает вычислительную стабильность. Блок \texttt{hidden\_4} выявлен как основной «двигатель» адаптации модели.

\vspace{0.8em}

\noindent \textbf{[Слайд 16: Заключение]} \\
Итог работы: метод PINN эффективен для уравнений высокого порядка. Понижение порядка нецелесообразно. Анализ NTK математически обосновал превосходство базовой архитектуры.

\vspace{1em}

\noindent Спасибо за внимание! Готова ответить на ваши вопросы.

\end{document}